\chapter{Landasan Kepustakaan}\label{bab:landasan}

\section{Kajian Pustaka}\label{sec:kajian-pustaka}

\subsection{Fusion Teks dan Citra pada Tobacco3482}

\textcite{audebert2020} merancang jaringan dua cabang yang dapat dilatih
\textit{end-to-end}:
\begin{itemize}
  \item cabang citra: MobileNetV2 dengan masukan $384\times384$;
  \item cabang teks: FastText dan CNN~1D atas teks hasil Tesseract~4;
  \item fusion: konkatenasi vektor fitur kedua cabang.
\end{itemize}
Pada Tobacco3482, fusion mencapai akurasi 87,8\%, di atas baseline citra (84,5\%) dan
baseline teks (73,8\%). Hasilnya dirata-ratakan atas tiga run tanpa simpangan baku.
Penulis juga melaporkan bahwa penjumlahan fitur berada di bawah baseline citra, tetapi
tanpa menyertakan angkanya. Di bagian penutup, penulis menyebut embedding yang tahan derau
OCR sebagai arah penelitian lanjutan. Paper inilah yang direplikasi pada studi
pendahuluan kami (Subbab~\ref{sec:studi-pendahuluan}).

\textcite{kashyap2026} adalah penerus langsung pendekatan tersebut dan menjadi
\textbf{referensi utama} penelitian ini. Ada empat perbedaan dari Audebert:
\begin{itemize}
  \item cabang citra memakai EfficientNetB3 dengan masukan $300\times300$;
  \item cabang teks memakai BERT-base;
  \item fusion dilakukan di tingkat keputusan, berupa rata-rata logit tanpa bobot;
  \item kedua cabang di-\textit{fine-tune} pada $\pm$320.000 dokumen RVL-CDIP, lalu
  Tobacco3482 dipakai hanya sebagai data uji.
\end{itemize}
Akurasi yang dilaporkan 93,06\% pada RVL-CDIP dan 92,00\% pada Tobacco3482.

Pembacaan teks lengkap paper tersebut menghasilkan tiga catatan:
\begin{enumerate}
  \item Mesin OCR-nya hanya disebut \enquote{mesin OCR komersial}, tanpa nama.
  \item Angka Tobacco3482 tidak konsisten dengan protokolnya sendiri. Kelas Note dan
  Report tidak ada di RVL-CDIP, sehingga akurasi maksimum pada 10 kelas hanya 86,6\%.
  Perhitungannya ada di Lampiran~\ref{lam:kashyap}.
  \item Hasil dilaporkan dari satu run tanpa uji statistik.
\end{enumerate}
Penulis sendiri menyebut pembobotan modalitas adaptif dan kondisi degradasi OCR yang
parah sebagai arah lanjutan. Keduanya bersinggungan langsung dengan penelitian ini.

\subsection{Gambaran Bidang: Meta-analisis Fusion untuk Klasifikasi Dokumen}

\textcite{mironczuk2026} menganalisis 139 studi primer fusion multimodal dan multiview
untuk klasifikasi dokumen. Fusion multimodal menaikkan akurasi rata-rata 5,28 poin
($p = 0{,}0016$). Namun efeknya terhadap F1 tidak signifikan secara statistik.

Tiga temuan lain dari review ini relevan langsung untuk penelitian ini:
\begin{enumerate}
  \item Hanya 11,8\% studi multimodal yang memakai uji statistik.
  \item Kualitas data, termasuk galat OCR, diakui sebagai sumber heterogenitas yang
  \emph{tidak dapat dikontrol} dalam meta-analisis, karena studi primer tidak
  melaporkannya.
  \item Penulis menganjurkan pelaporan baseline unimodal, F1 untuk data timpang, dan
  ablasi lengkap untuk setiap modalitas yang ditambahkan.
\end{enumerate}

\subsection{Perbandingan Mesin OCR dan Dampaknya pada Tugas Lanjutan}

\textcite{francis2025} membandingkan tujuh mesin OCR, termasuk Tesseract, EasyOCR,
PaddleOCR, dan Keras-OCR, pada citra campuran: dokumen cetak, layar, struk, dan ekspresi
matematika. PaddleOCR mencatat akurasi tertinggi (0,7136), diikuti EasyOCR (0,6953),
Tesseract (0,6868), dan Keras-OCR (0,6784). Selisihnya kecil, sekitar 3,5 poin. Yang
diukur di sana adalah ketepatan OCR itu sendiri, bukan dampaknya pada tugas yang
memakai teks hasil OCR.

Dampak semacam itu diukur oleh \textcite{zhang2025} untuk sistem
\textit{retrieval-augmented generation}. Mereka menyusun tolok ukur OHRBench dari 8.561
citra dokumen. Galat OCR dipilah menjadi derau semantik (kata salah baca) dan derau
format (struktur tabel atau rumus rusak). Keduanya terbukti menurunkan kinerja tugas
lanjutan. Untuk klasifikasi dokumen, pengukuran serupa belum kami temukan.

\subsection{Ketahanan Embedding terhadap Derau OCR}

\textcite{michail2025} menunjukkan bahwa model embedding standar melemah pada teks hasil
OCR, karena derau OCR jarang muncul dalam data pelatihannya. Ketahanannya dapat
dipulihkan dengan cara yang murah: \textit{fine-tuning} memakai teks yang diberi derau
karakter buatan dan \textit{loss} kontrastif, tanpa data berlabel dari domain sasaran.
Temuan ini mendasari sumbu embedding dalam penelitian ini. FastText menangani derau OCR
secara struktural melalui \textit{subword}, sedangkan BERT menanganinya melalui
tokenisasi WordPiece dan konteks kalimat.

\subsection{Jalan Pintas pada Tobacco3482}

\textcite{larson2025} menunjukkan bahwa dokumen RVL-CDIP dan Tobacco3482 bercap
\textbf{kode ID} arsip. Model sederhana yang hanya memakai fitur dari kode ID itu sudah
mencapai $\pm$60\% akurasi di Tobacco3482. Pada RVL-CDIP, classifier mutakhir kehilangan
11 poin akurasi setelah kode ID dihapus.

Kode ID ini terlihat pada contoh dokumen di Gambar~\ref{fig:contoh-pasangan}, misalnya
\texttt{00625137} dan \texttt{CTR CONTRACTS 023247}. Temuan ini menjadi kontrol wajib
pada tahap lanjutan yang memperlakukan teks sebagai objek berposisi
(Subbab~\ref{sec:lanjutan}).

\subsection{Posisi Penelitian}

Tabel~\ref{tab:posisi} memetakan studi-studi di atas terhadap lima aspek yang menjadi
fokus penelitian ini.

\begin{table}[H]
  \centering
  \caption{Posisi penelitian terhadap studi terkait}
  \label{tab:posisi}
  \fontsize{10}{12}\selectfont
  \begin{tabularx}{\textwidth}{|>{\raggedright\arraybackslash}X|c|c|c|c|c|}
    \hline
    \textbf{Studi} & \textbf{\makecell{Mesin OCR\\disebut}} &
    \textbf{\makecell{$>1$\\mesin OCR}} & \textbf{\makecell{Variasi\\embedding}} &
    \textbf{\makecell{Uji missing\\modality}} & \textbf{\makecell{Uji\\statistik}} \\ \hline
    \textcite{audebert2020} & \checkmark & -- & MLP vs CNN & -- & -- \\ \hline
    \textcite{kashyap2026} & -- & -- & -- & -- & -- \\ \hline
    \textcite{francis2025} & \checkmark & \checkmark\ (7) &
    \multicolumn{3}{c|}{tanpa tugas klasifikasi} \\ \hline
    \textbf{Penelitian ini} & \checkmark & \checkmark\ (3) & FastText, BERT &
    \checkmark & \checkmark \\ \hline
  \end{tabularx}
  \sumber{Dirangkum dari studi yang dirujuk.}
\end{table}

\section{Dasar Teori}\label{sec:dasar-teori}

\subsection{Optical Character Recognition}

OCR modern umumnya bekerja dalam dua tahap: \textbf{deteksi} wilayah teks pada citra,
lalu \textbf{pengenalan} karakter di dalam setiap wilayah. Ketiga mesin yang diuji
berbeda pada kedua tahap tersebut.

\textbf{Tesseract} \parencite{smith2007} memulai dengan analisis tata letak halaman untuk
menemukan blok, baris, dan kata, lalu mengenali teksnya. Sejak versi 4, pengenalannya
memakai jaringan LSTM yang membaca setiap baris sebagai sekuens. Konfigurasi
\texttt{-{}-psm 3} berarti segmentasi halaman otomatis penuh tanpa deteksi orientasi.
Konfigurasi inilah yang dipakai QS-OCR \parencite{qsocr}. Keluarannya mengikuti urutan baca
Tesseract sendiri, yaitu blok, paragraf, baris, lalu kata.

\textbf{EasyOCR} \parencite{easyocr} mendeteksi teks dengan CRAFT \parencite{baek2019}. CRAFT
memprediksi wilayah setiap karakter dan keterkaitan antar-karakter, lalu menggabungkannya
menjadi kotak kata atau baris. Pengenalannya memakai CRNN \parencite{shi2017}, yang terdiri
atas tiga bagian:
\begin{itemize}
  \item ekstraktor fitur konvolusi;
  \item LSTM dua arah untuk pelabelan sekuens;
  \item dekoder \textit{connectionist temporal classification} (CTC).
\end{itemize}

\textbf{PaddleOCR} \parencite{paddleocr} menyusun pipeline dari modul-modul terpisah:
klasifikasi orientasi dokumen, koreksi distorsi, deteksi baris teks, klasifikasi
orientasi baris, dan pengenalan teks. Modul orientasi dan koreksi distorsi bersifat
opsional. Penelitian ini memakai pasangan model deteksi dan pengenalan PP-OCRv5
\textit{mobile}.

Mesin berbasis kotak seperti EasyOCR dan PaddleOCR tidak menjamin urutan baca. Urutan
itu harus direkonstruksi dari posisi kotak (Subbab~\ref{sec:metode-ocr}).

\subsection{Representasi Kata}

\textbf{FastText} \parencite{bojanowski2017} merepresentasikan kata $w$ sebagai jumlah
vektor seluruh \textit{n-gram} karakternya:
\begin{equation}\label{eq:fasttext}
  \mathbf{u}_w = \sum_{g \in \mathcal{G}_w} \mathbf{z}_g
\end{equation}
dengan $\mathcal{G}_w$ himpunan \textit{n-gram} karakter kata $w$ beserta kata itu
sendiri, dan $\mathbf{z}_g$ vektor untuk \textit{n-gram} $g$. Karena itu, kata yang rusak
akibat OCR, misalnya \texttt{Tospacco}, tetap mendapat vektor yang dekat dengan kata
aslinya, karena keduanya berbagi banyak \textit{n-gram}.

\textbf{BERT} \parencite{devlin2019} adalah encoder Transformer dua arah yang dilatih awal
dengan tugas \textit{masked language modelling}. Teks dipecah menjadi sub-kata dengan
tokenisasi WordPiece, dan vektor setiap token bergantung pada konteks kalimatnya.
Panjang masukan dibatasi 512 token. Untuk klasifikasi, representasi token khusus
\texttt{[CLS]} dilewatkan ke lapisan klasifikasi, dan seluruh model di-\textit{fine-tune}.

\subsection{Arsitektur Konvolusi}

\textbf{CNN~1D untuk teks} \parencite{kim2014} menggeser filter konvolusi di sepanjang
sekuens vektor kata, lalu meringkasnya dengan \textit{max-pooling} terhadap waktu. Filter
dengan lebar $k$ berfungsi sebagai pendeteksi frasa sepanjang $k$ kata. Model ini peka
terhadap \textbf{urutan} kata, berbeda dari rata-rata vektor kata.

\textbf{MobileNetV2} \parencite{sandler2018} adalah CNN citra berbiaya rendah. Blok
penyusunnya adalah \textit{inverted residual} dengan \textit{linear bottleneck}: kanal
diekspansi dengan konvolusi $1\times1$, difilter dengan konvolusi \textit{depthwise}
$3\times3$, lalu diproyeksikan kembali tanpa aktivasi non-linear.

\subsection{Strategi Fusion}

Fusion dapat dilakukan pada tingkat fitur atau tingkat keputusan
\parencite{baltrusaitis2019}. Misalkan $\mathbf{v}_t, \mathbf{v}_v \in \mathbb{R}^{d}$
adalah vektor fitur cabang teks dan citra, dan $\mathbf{z}_t, \mathbf{z}_v \in
\mathbb{R}^{C}$ adalah logit keduanya untuk $C$ kelas.

\textbf{Fusion tingkat fitur} menggabungkan vektor fitur lalu mengklasifikasikannya
bersama. Pada konkatenasi \parencite{audebert2020}:
\begin{equation}\label{eq:concat}
  \hat{y} = \arg\max_c \; \mathrm{MLP}\!\left([\,\mathbf{v}_t \,;\, \mathbf{v}_v\,]\right)_c
\end{equation}

\textbf{Fusion tingkat keputusan} menggabungkan keluaran setiap cabang. Pada rata-rata
logit tanpa bobot \parencite{kashyap2026}:
\begin{equation}\label{eq:logit}
  \hat{y} = \arg\max_c \; \tfrac{1}{2}\left(z_{t,c} + z_{v,c}\right)
\end{equation}
Fusion tingkat keputusan menjamin kedua cabang ikut menentukan hasil. Akibatnya, cabang
yang lemah tetap memengaruhi keputusan.

Pada fusion tingkat fitur, cabang dengan skala fitur jauh lebih besar dapat mendominasi.
Masalah ini dapat diatasi dengan \textbf{Layer Normalization} \parencite{ba2016}, yang
menormalkan setiap vektor fitur terhadap rata-rata $\mu$ dan simpangan baku $\sigma$
komponennya sendiri:
\begin{equation}\label{eq:layernorm}
  \mathrm{LN}(\mathbf{v}) = \boldsymbol{\gamma} \odot \frac{\mathbf{v} - \mu}{\sigma +
  \epsilon} + \boldsymbol{\beta}
\end{equation}
dengan $\boldsymbol{\gamma}$ dan $\boldsymbol{\beta}$ parameter yang dipelajari. LN
bekerja per sampel, sehingga tidak bergantung pada ukuran \textit{batch}.

\subsection{Metrik Evaluasi}

\textbf{Overall accuracy} (OA) adalah proporsi dokumen uji yang diklasifikasikan dengan
benar. Kelas Tobacco3482 timpang, dengan rasio kelas terbesar berbanding terkecil
5,2:1, sehingga OA cenderung didominasi kelas besar. Karena itu \textbf{macro F1} juga
dilaporkan:
\begin{equation}\label{eq:macrof1}
  \mathrm{F1}_c = \frac{2\, P_c R_c}{P_c + R_c}, \qquad
  \text{macro F1} = \frac{1}{C} \sum_{c=1}^{C} \mathrm{F1}_c
\end{equation}
dengan $P_c$ dan $R_c$ presisi dan \textit{recall} kelas $c$.

\textbf{Oracle} adalah batas atas teoretis fusion: sebuah dokumen dihitung benar bila
cabang teks \emph{atau} cabang citra mengklasifikasikannya dengan benar
\parencite{audebert2020}. Selisih antara oracle dan baseline terbaik adalah potensi yang
dapat dipanen fusion.

\textbf{Uji \textit{missing modality}} mengukur ketergantungan model gabungan pada setiap
modalitas. Caranya, masukan satu modalitas dinolkan saat inferensi:
\begin{equation}\label{eq:delta}
  \Delta_\text{teks} = \mathrm{OA}_\text{utuh} - \mathrm{OA}_\text{teks dinolkan}
\end{equation}
$\Delta_\text{citra}$ didefinisikan secara analog. $\Delta_\text{teks}$ yang mendekati nol
berarti model tidak memakai teks, berapa pun akurasinya.

\textbf{Uji McNemar} \parencite{mcnemar1947} membandingkan dua classifier pada data uji
yang sama. Uji ini dianjurkan untuk perbandingan algoritme pembelajaran oleh
\textcite{dietterich1998}. Bila $b$ adalah jumlah dokumen yang hanya benar pada classifier
pertama dan $c$ jumlah dokumen yang hanya benar pada classifier kedua:
\begin{equation}\label{eq:mcnemar}
  \chi^2 = \frac{(|b - c| - 1)^2}{b + c}
\end{equation}
Statistik ini berdistribusi $\chi^2$ dengan satu derajat bebas. Bila $b + c$ kecil,
dipakai uji binomial eksak. Untuk $m$ perbandingan sekaligus, ambang signifikansi
dikoreksi menjadi $\alpha/m$ (koreksi Bonferroni, \cite{dunn1961}).
